% appx-lockstep.tex — Appendix C: Python <-> Lean lockstep test
% enumeration.

\chapter{Python$\leftrightarrow$Lean Lockstep Tests}
\label{ch:appx-lockstep}

\noindent\emph{The test surface that binds the Python numerical
pilot to the Lean rational literals, ensuring that every constant
the Lean witness depends on agrees with the Python computation to
the last $\Qp$-rational digit.}

The manuscript's two-layer construction --- $\Qp$-rational data
theorems in Lean (\Cref{ch:appx-lean-audit}) on top of Python-side
numerical pilots (\Cref{ch:pilot-su2}) --- requires that the
numerical values in the Python pilots match the rational literals
hardcoded in the Lean library exactly.  This appendix enumerates
the tests that enforce this lockstep: a drift in either side
without the corresponding update on the other triggers a CI
failure.

\section{The lockstep contract}
\label{sec:lockstep-contract}

\marginkey{Two-layer contract}%

The Python pilot computes per-step $\etak$ values as floats; the
Lean witness uses $\Qp$-rational literals with explicit
denominators.  The contract is:

\begin{enumerate}[label=(L\arabic*)]
  \item Every Lean $\Qp$-literal that appears in a load-bearing
    theorem has a corresponding Python computation that produces
    the same numerator at the matching denominator (with the
    Python value $\le$ the Lean literal at every operating point,
    so the Lean witness is an honest upper bound).
  \item The lockstep tests in \codepath{tests/test\_q\_lift.py}
    assert numerator equality (or "Python $\le$ Lean rational")
    for each tracked literal.
  \item A Python-side numerical change that violates the contract
    triggers a test failure; updating the contract requires
    deliberate co-edits to both the Python ceiling and the Lean
    literal.
\end{enumerate}

\section{Tracked tests}
\label{sec:lockstep-tracked}

\begin{table}[h]
\centering
\renewcommand{\arraystretch}{1.18}
\begin{tabular}{p{4.5cm}@{\hspace{1em}}p{4.5cm}@{\hspace{1em}}p{4cm}}
\toprule
Python test                       & Asserts                          & Lean literal                \\
\midrule
\codepath{test\_eta\_max\_matches\_lean\_hardcoded}
& Max per-step $\etak$ on the canonical pilot ($K = 3$) is
  bounded by the Lean rational.
& \codepath{eta\_max\_pilot}     \\
\codepath{test\_eta\_total\_matches\_lean\_hardcoded}
& Cumulative $\sum_{k < 3} \etak$ canonical bounded by the Lean
  literal.
& \codepath{eta\_total\_pilot}   \\
\codepath{test\_wilson\_eta\_max\_matches\_lean\_hardcoded}
& Same for the Wilson pilot at $g_0 = 0.50$, $K = 3$.
& \codepath{eta\_max\_pilot\_wilson} \\
\codepath{test\_wilson\_eta\_total\_matches\_lean\_hardcoded}
& Cumulative Wilson-pilot sum vs.\ Lean literal.
& \codepath{eta\_total\_pilot\_wilson} \\
\codepath{test\_canonical\_pilot\_data\_matches\_lean\_module}
& Full per-scale $\etak$ vector at $K = 3$ matches the Lean module's
  per-scale rational table.
& \codepath{RG.SigmaHatPilot}    \\
\codepath{test\_canonical\_pilot\_closes}
& The five-criterion validator predicate
  $\mathrm{Pilot}_{\SUtwo}(0.30, 3) = \mathsf{True}$ holds in
  Python, mirroring the Lean
  \codepath{eta\_pilot\_budget\_closes}.
& \codepath{eta\_pilot\_budget\_closes} \\
\bottomrule
\end{tabular}
\caption{Python$\leftrightarrow$Lean lockstep tests in
\codepath{tests/test\_q\_lift.py} and
\codepath{tests/test\_pilot\_su2.py}.  Each row binds a Python
numerical assertion to a Lean rational literal; CI failure
indicates drift in one or the other.}
\label{tab:lockstep-tests}
\end{table}

\section{The rational-ceiling helper}
\label{sec:rational-ceiling}

\marginkey{Float $\to$ $\Qp$ ceiling}%

The bridge from a Python float to a $\Qp$-rational upper bound is
the helper \codepath{rational\_ceiling(x, denom)} in
\codepath{computations/q\_lift.py}: it returns the smallest integer
$n$ with $n/\mathrm{denom} \ge x$, i.e., the rational ceiling at
the given denominator.  All Lean rationals in the load-bearing
theorems are sourced from this helper applied to the Python
numerical output, ensuring the Lean literal is always a
\emph{conservative} (upper-bound) rational, never an underestimate
of the true Python value.

The properties asserted in the test suite:
\begin{itemize}[nosep]
  \item \codepath{test\_exact\_rational}:
    $\mathrm{ceiling}(p/q, q) = p$ when $x = p/q$ is exact.
  \item \codepath{test\_strict\_round\_up}:
    irrational $x$ rounds strictly up.
  \item \codepath{test\_upper\_bound\_property}:
    $\mathrm{ceiling}(x, \mathrm{denom}) / \mathrm{denom} \ge x$
    for every valid input.
  \item \codepath{test\_invalid\_denominator}: rejects $\mathrm{denom} \le 0$.
\end{itemize}

\section{Reproduction}
\label{sec:lockstep-reproduce}

The lockstep tests are part of the project's main
\codepath{pytest} suite.  From the repository root:
\begin{itemize}
  \item \codepath{pytest tests/test\_q\_lift.py -v}
    runs the Python$\leftrightarrow$Lean rational tests in
    isolation; the manuscript-bundle release pipeline runs
    these as part of the broader $500+$-test suite.
  \item \codepath{pytest -k lockstep} runs every test whose name
    contains ``lockstep'', spanning multiple test files.
\end{itemize}

The combined Python$\leftrightarrow$Lean tests pass in under five
seconds on a modern laptop and have been a stable manuscript-build
gate since their introduction.
